\section{Results and discussion}
\label{sec:results}

% Source of every agentic and trained-model number in this section unless noted:
% outputs/run_20261109_121625_bench_agentic_arms_beside_the_20M_legacy/scorecard.md
% (houses_32.s0.70-15-15.v1, seed 0, reference 399, 17 arms, permutation n_perm=10000, Holm).
% p_Holm = 0.0136 is the smallest Holm-adjusted p the card reports; we print it as 0.014.

\begin{table*}[t]
\centering
\small
\setlength{\tabcolsep}{4pt}
\caption{\textbf{Agentic arms on the house benchmark, $n{=}133$ test prompts.} MV-KID reads appearance and Geom-KID geometry (lower is better; letters: arms sharing a letter are not separated after Holm correction). Enterable is the share of builds a player walks into (real reference 0.945). Twin is the share of builds with a near-copy in the same set. Trained models are scored on the same card at $n{=}256$, so their BlockScore is not comparable with an $n{=}133$ arm (\cref{sec:setup}). Brackets are 95\% intervals.}
\label{tab:main}
\begin{tabular}{@{}llcccccl r@{}}
\toprule
Arm & $n$ & MV-KID $\downarrow$ & Geom-KID $\downarrow$ & Enterable & Twin & Diversity & Gate & \$/build \\
\midrule
Real (held out) & 133 & $-0.001$\,\textsuperscript{a} & $-0.08$\,\textsuperscript{a} & 0.932 & 0.000 & 0.902 & pass & -- \\ % scorecard real_test rows
\midrule
\textsc{One-Shot} & 133 & \multicolumn{6}{c}{\TODO{bench One-Shot}} & 0.010 \\ % run_20261109_130512_bench_houses32_oneshot_test: paid $1.3892 / 133 builds; not yet benched
\textsc{Classic-Matched} & 133 & 0.403\,\textsuperscript{ef} & 18.35\,\textsuperscript{ef} & 0.677 & 0.316 & 0.833 & DQ twin & 0.045 \\ % scorecard realism/traversal/novelty/cost/leaderboard
\textsc{Dream} & 133 & 0.342\,\textsuperscript{e} & 17.88\,\textsuperscript{ef} & 0.496 & 0.301 & 0.786 & DQ twin & 0.040 \\ % scorecard
\textsc{Dream-Caption} & 133 & \textbf{0.217}\,\textsuperscript{d} & \textbf{15.71}\,\textsuperscript{de} & 0.511 & \textbf{0.211} & 0.790 & DQ twin & 0.048 \\ % scorecard
\midrule
\multicolumn{9}{@{}l}{\emph{Trained voxel models (reference rows)}} \\
\texttt{canon\_oriented} 5M & 256 & 0.254\,\textsuperscript{d} & 14.24\,\textsuperscript{d} & 0.605 & 0.000 & 0.885 & pass (150.5) & -- \\ % scorecard canon_oriented rows; BlockScore 150.53
20M legacy & 256 & 0.159\,\textsuperscript{c} & 9.06\,\textsuperscript{c} & 0.629 & 0.008 & 0.878 & DQ mem. & -- \\ % scorecard recanon20m_s0 rows; T118 model
\bottomrule
\end{tabular}
\end{table*}
% Intervals for the agentic rows (scorecard realism table):
% MV-KID: classic 0.403 [0.361, 0.450], dream 0.342 [0.311, 0.375], caption 0.217 [0.200, 0.237]
% Geom-KID: classic 18.351 [17.056, 19.534], dream 17.882 [17.028, 18.980], caption 15.712 [14.888, 16.681]
% Twin: classic 0.316 [0.241, 0.391], dream 0.301 [0.226, 0.383], caption 0.211 [0.143, 0.286]

%TL;DR We organize the results as six findings, each tied to one row group of the main table; the dream buys appearance through its caption and buys nothing in geometry or structure.
\Cref{tab:main} holds the main comparison. Each subsection below states one finding and the test behind it.

\subsection{The caption, not the image, moves appearance}

%TL;DR Dream-Caption is the only agentic arm separated from the compute-matched control on MV-KID; Dream alone is not.
\textsc{Dream-Caption} reaches MV-KID 0.217, against 0.342 for \textsc{Dream} and 0.403 for \textsc{Classic-Matched}. % scorecard realism
It beats \textsc{Dream} by 0.125 (95\% interval $[0.089, 0.164]$, $p_{\text{Holm}}{=}0.014$) and \textsc{Classic-Matched} by 0.186 ($[0.135, 0.229]$, $p_{\text{Holm}}{=}0.014$). % scorecard mv_dino_kid h2h: dream-caption vs dream -0.125, vs classic-matched -0.186
\textsc{Dream} and \textsc{Classic-Matched} share a letter, so the image alone does not separate from the compute-matched control. % scorecard groups e / ef
Because the two dream arms build from identical cached images, the 0.125 gap belongs to the caption. % dream-caption.yaml shared cache
\textsc{Dream-Caption} also shares its letter with the trained \texttt{canon\_oriented} model (0.254), so on appearance a prompted agent with a captioned image is not separated from a 5M model trained on the corpus. % scorecard groups d / d

%TL;DR Small pilots show what the caption does: it names parts the builder misreads in the image and lists scenery to leave out.
Pilots show the mechanism, at sizes too small to test. % T46 n=4; T48 n=6
In a four-prompt pilot, the caption arm was the only arm whose windmill ended with upright sails, after the caption named ``four large diagonal sail arms''. % T46
In a six-prompt local pilot with SDXL-Turbo dreams that place every house on terrain, the caption listed the terrain as surroundings to omit and the builds dropped it (coherence rate 1.00 against 0.83 without the caption). % T48: coherence 1.00 vs 0.83
In the same four-prompt pilot, a caption from a different prompt turned the lighthouse into a red box house. % T46 dream-caption-shuffled
\textsc{Dream-Caption-Shuffled} tests this at full $n$ (\cref{tab:controls}).

%TL;DR The caption arm builds smaller houses than real ones, which is part of how its renders differ.
The caption also shrinks the builds: \textsc{Dream-Caption} places 704 blocks with a largest extent of 13.4, against 953 and 15.0 for \textsc{Dream}, 1{,}202 and 19.1 for \textsc{Classic-Matched}, and 904 and 19.2 for real houses. % scorecard coherence table: n_blocks, max_dim
\TODO{check whether MV-KID renders normalize build size; if not, measure how much of the caption's MV-KID gain survives a size-matched comparison.}

\subsection{Geometry: the three agentic arms tie}

%TL;DR No pair of agentic arms separates on Geom-KID, so neither the image nor the caption moves the distribution of shapes.
On Geom-KID the three arms score 15.71, 17.88 and 18.35, and every pair shares a letter. % scorecard geom_kid groups de / ef / ef
The image and the caption leave the shape distribution where the compute-matched control puts it.
\textsc{Dream} and \textsc{Classic-Matched} also share a letter with the protocol's Patchwork baseline (18.79), and only \textsc{Dream-Caption} separates from it ($-3.08$, $[-4.52, -1.55]$, $p_{\text{Holm}}{=}0.014$). % scorecard: patchwork 18.791 group f; patchwork vs dream-caption +3.080
All three beat the solidified-real control (21.56), whose interiors are filled. % scorecard: real@solidify 21.562 group g; vs classic-matched p_holm 0.0425

\subsection{Structure: the dream arms are harder to enter}

%TL;DR A player walks into 68% of Classic-Matched houses but only half of the dream houses, against 93% of real ones.
The traversal probe orders the arms the other way from MV-KID. % scorecard traversal table
A player enters 0.677 of \textsc{Classic-Matched} builds, 0.511 of \textsc{Dream-Caption} builds and 0.496 of \textsc{Dream} builds, against 0.932 for held-out real houses. % scorecard enterable
In real-sample spreads the distance to real is 1.17 for the control and 1.90 and 1.96 for the dream arms. % scorecard traversal: enterable spreads
Players reach 0.77 storeys per \textsc{Classic-Matched} build and 0.51 to 0.53 per dream build, against 1.57 for real houses. % scorecard floors_reached
The protocol reports traversal as a distance to real without a pairwise test. \TODO{run a paired permutation test on per-build enterability between Classic-Matched and each dream arm.}

\begin{table}[t]
\centering
\footnotesize
\setlength{\tabcolsep}{3pt}
\caption{\textbf{Inspection failures.} Share of builds failing each text geometry check at warning severity or worse; checks with no failures in any row are omitted. Real houses fail mostly through crop artefacts (missing floors, detached yard pieces).}
\label{tab:inspect}
\begin{tabular}{@{}lrrrrr@{}}
\toprule
Arm & $n$ & Clean & Roofed & Enterable & Headroom \\
\midrule
\textsc{Classic-Matched} & 133 & 68\% & 6\% & 33\% & 1\% \\
\textsc{Dream} & 133 & 53\% & 2\% & 54\% & 7\% \\
\textsc{Dream-Caption} & 133 & 48\% & 0\% & 57\% & 13\% \\
Real houses & 400 & 55\% & 10\% & 17\% & 9\% \\
\bottomrule
\end{tabular}
\end{table}
% Source: T95 (blockgen/agentic/inspection.py on the T51 builds and 400 real houses_32 builds, seed 1).
% Real also fails connected 6%, supported 14%; agentic 0% on both; dream-caption interior 2%.

%TL;DR An independent rule-based inspector agrees: the enterable check fails on more than half of the dream builds and a third of the control builds.
A rule-based inspector, run on the same builds, agrees (\cref{tab:inspect}). % T95
The enterable check fails on 54\% of \textsc{Dream} builds and 57\% of \textsc{Dream-Caption} builds, against 33\% for \textsc{Classic-Matched} and 17\% for real houses. % T95
Headroom failures rise from 1\% for the control to 7\% and 13\% for the dream arms. % T95
The agentic builds never fail the connected or supported checks, which real crops fail at 6\% and 14\%. % T95
The dream arms build connected, supported shells that a player does not get into.

%TL;DR All agentic arms leave interiors too small and too symmetric, and the image moves neither.
Interiors are the shared gap. % scorecard coherence / geometry scalars
The enclosed-air ratio is 0.039 to 0.043 for the three arms against 0.114 for real houses, and yaw symmetry is 0.92 to 0.96 against 0.68. % scorecard: enclosed_air_ratio, yaw_symmetry
The target is an exterior view: it fixes facades, roof and materials and says nothing about rooms, doors or stairs. % dream.yaml: dreams [iso]
The appearance gain and the structural loss are consistent with a builder that spends its edits on what the critic sees.
\TODO{test the interpretation with a cutaway or plan dream at n=256 (recipes dream-multiview, blueprint exist; T42 found the plan dream redraws the exterior).}

\subsection{Template collapse: every agentic arm fails the near-twin gate}

%TL;DR A fifth to a third of each agentic set has a near-copy inside the same set, against none of the real houses, so all three arms are disqualified.
Within each set, 31.6\% of \textsc{Classic-Matched} builds, 30.1\% of \textsc{Dream} builds and 21.1\% of \textsc{Dream-Caption} builds have a near-twin, against 0.0\% of real houses. % scorecard novelty: voxel_twin_rate
The gate disqualifies all three, at 15.8, 15.0 and 10.5 real spreads above the real rate. % scorecard leaderboard DQ reasons
Diversity agrees: 0.833, 0.786 and 0.790 against 0.902 for real houses. % scorecard voxel_diversity
On inspection, \textsc{Classic-Matched} draws most prompts as the same oak cottage with a gabled roof. % T51 (by eye)
The caption lowers the twin rate by 0.105 against the control, inside overlapping intervals ($[0.143, 0.286]$ against $[0.241, 0.391]$). % scorecard
Both trained models pass the near-twin gate (0.000 and 0.008). % scorecard

%TL;DR Before the twin gate existed, BlockScore ranked Dream-Caption above the other agentic arms and an LLM baseline, which shows why the gate is required.
Before the near-twin gate entered the protocol, BlockScore ranked \textsc{Dream-Caption} (155.8) above \textsc{Dream} (177.2) and \textsc{Classic-Matched} (181.9), and above an LLM block-serialization baseline (167.9). % T64 (n=128 card, 2026-10-23), before T91 (2026-11-04)
A leaderboard without the gate rewards an agent that repeats one house. % T91: ungated vanilla templates 202.95 / 194.98 ahead of every BlockGen model
No memorization gate fires on an agentic arm: DINO-NN percentiles are 0.460 to 0.591 against 0.536 for real houses. % scorecard dino_nn_percentile

\subsection{Trained models still lead}

%TL;DR The 20M trained model beats the best agentic arm on both co-primary metrics, and the 5M model beats Dream and Classic-Matched on geometry.
The 20M legacy voxel model beats \textsc{Dream-Caption} on Geom-KID by 6.65 ($[5.24, 8.08]$, $p_{\text{Holm}}{=}0.014$) and on MV-KID by 0.058 ($[0.030, 0.083]$, $p_{\text{Holm}}{=}0.014$). % scorecard h2h: recanon20m_s0 vs dream-caption -6.649, -0.058
The 5M \texttt{canon\_oriented} model beats \textsc{Dream} by 3.64 and \textsc{Classic-Matched} by 4.11 on Geom-KID (both $p_{\text{Holm}}{=}0.014$) and ties \textsc{Dream-Caption} on both metrics. % scorecard h2h: canon_oriented vs dream -3.638, vs classic -4.106; groups d/de, d/d
\texttt{canon\_oriented} is the only arm in \cref{tab:main} with a BlockScore (150.5). % scorecard leaderboard rank 2
On enterability \textsc{Classic-Matched} (0.677) exceeds both trained models (0.605 and 0.629). % scorecard traversal

%TL;DR The 20M model fails the memorization gate on this card and passes it on the n=256 card, which is the equal-n rule at work.
The 20M model fails the memorization gate here (DINO-NN percentile 0.359 against a real 0.536, 3.5 spreads) and passes every gate on an $n{=}256$ card. % scorecard leaderboard; T118 ("No gate fires on a model arm")
The gate's spread comes from a real sample of the card's $n$, so this verdict is one more reason to read the agentic and trained rows only through the pairwise tests. % docs/benchmark.md Protocols
The trained arms are scored at $n{=}256$ and the agentic arms at $n{=}133$; KID is unbiased in $n$, but the tests mix sample sizes. % scorecard n column
\TODO{n=256: rerun this comparison with all arms at n=256.}

\subsection{Cost}

%TL;DR A loop build costs 3.8 to 4.6 times a one-shot build, and the dream arms cost within 11% of the compute-matched control.
The three loop arms cost \$0.040 to \$0.048 per build with 3.5 to 4.6 model calls, against \$0.010 and one call for \textsc{One-Shot}, a factor of 3.8 to 4.6. % scorecard cost; run summaries n_llm_calls; oneshot paid $1.3892/133 = 0.01045; 0.040/0.01045=3.8, 0.048/0.01045=4.6
\textsc{Dream} costs 11\% less than \textsc{Classic-Matched} and \textsc{Dream-Caption} 7\% more. % 0.040/0.045, 6.385/5.968
Completion tokens per build are 14{,}600 for \textsc{Dream}, 17{,}962 for \textsc{Dream-Caption} and 20{,}315 for \textsc{Classic-Matched}, against 4{,}958 for \textsc{One-Shot}. % scorecard cost completion_tokens; oneshot summary.md
At these prices the full $n{=}133$ run of an arm costs \$5.30 to \$6.39. % run summaries total_would_cost_usd: 5.299, 5.968, 6.385
The trained models cost no API calls at sampling time. \TODO{add GPU-hours for training canon\_oriented and the 20M legacy model (T111, T118) so the cost comparison is complete.}

\subsection{Controls and the $n{=}256$ run}

\begin{table}[t]
\centering
\footnotesize
\setlength{\tabcolsep}{3pt}
\caption{\textbf{Controls and the $n{=}256$ card.} Each row isolates one component. All rows use the 256 prompts of \cref{sec:setup}.}
\label{tab:controls}
\resizebox{\columnwidth}{!}{%
\begin{tabular}{@{}llcc@{}}
\toprule
Arm & Isolates & MV-KID & Geom-KID \\
\midrule
\textsc{One-Shot} & the loop & \UNMEASURED{} & \UNMEASURED{} \\
\textsc{Classic-Matched} & (control) & \UNMEASURED{} & \UNMEASURED{} \\
\textsc{Dream} & the image & \UNMEASURED{} & \UNMEASURED{} \\
\textsc{Dream-Caption} & the caption & \UNMEASURED{} & \UNMEASURED{} \\
\textsc{Dream-Caption-Shuffled} & correct caption & \UNMEASURED{} & \UNMEASURED{} \\
\textsc{Text-Spec} & the picture & \UNMEASURED{} & \UNMEASURED{} \\
\textsc{Best-of-4} & more samples & \UNMEASURED{} & \UNMEASURED{} \\
\texttt{inspect-lessons} & text feedback & \UNMEASURED{} & \UNMEASURED{} \\
Gemma 4 12B (local) & the model & \UNMEASURED{} & \UNMEASURED{} \\
\bottomrule
\end{tabular}}
\end{table}
% State of outputs/logs/agentic256.log at 2026-11-09 20:58 UTC (in progress, not waited for):
%   One-Shot test (133) done: run_20261109_130512_bench_houses32_oneshot_test, paid $1.3892, not yet benched.
%   One-Shot val (123) done: run_20261109_145342_bench_houses32_oneshot_val, paid $1.3107, not yet benched.
%   Agentic val (123 x [dream-caption, dream, classic-matched]): run_20261109_165950_bench_houses32_agentic_val,
%   dream-caption at 86/123; dream and classic-matched val not started.
%   The test-split builds of the three loop arms already exist (T51 runs, dream-caption133).

%TL;DR The control table fixes in advance which comparison answers which question, so the n=256 card reads off directly.
\Cref{tab:controls} fixes the comparisons the $n{=}256$ card answers. \TODO{n=256}
\textsc{Dream-Caption} against \textsc{Dream-Caption-Shuffled} measures what a correct description adds beyond more text; if they tie, the gain is tokens. % docs/agentic-suite.md 5.6
\textsc{Dream-Caption} against \textsc{Text-Spec} measures what the picture adds beyond a description.
\textsc{Dream} against \textsc{Best-of-4} at matched cost measures whether the image beats sampling four times. % best-of-4.yaml comment; docs 9
\texttt{inspect-lessons} against \textsc{Classic-Matched} measures whether text geometry feedback fixes the enterability gap of \cref{tab:inspect}, which a render critic does not see. % inspection.py calibrated on 400 real houses (mean 0.93), memory note agentic-suite
The local Gemma arm measures whether the findings depend on the closed model. % T47, T48 pilots n=6

\subsection{Limitations}

%TL;DR One run per prompt, one builder, one image model, one building type and one canvas size per arm bound every claim above.
Every arm runs once per prompt with one builder (\texttt{gpt-5-mini}) and one image model (\texttt{gpt-image-1-mini}), so the findings carry no seed variance and no model variance. % summaries: n_builds = prompts; dream.yaml
The prompts are houses at $32^3$, and the $n{=}133$ card sits just above the protocol's minimum of 128 per arm. % docs/benchmark.md min_n 128
\textsc{One-Shot} ran on a $48^3$ canvas, unlike the loop arms. % oneshot config.json
The selectors differ between arms (CLIP text for the control, SigLIP image similarity for the dream arms), so a dream arm changes both its input and its selection signal. % recipes score fields
The titles name real GrabCraft builds that the builder may have seen during pretraining; the memorization gate checks proximity to our training split, not to the web. % scorecard dino_nn_percentile is against the training split
The $n{=}256$ prompt set adds validation titles while the reference stays the test split. % prompts files; scorecard reference 399
